% =====================================================================
% Section 1: Introduction — EvidenceFlow
% =====================================================================
%
% Status: Draft §1 Introduction v2 (2026-05-17)
% Revision note:
%   - Rewrite starts directly from GraphRAG, following advisor feedback.
%   - Narrative: GraphRAG routes -> limitations -> motivation -> method.
%   - Contribution bullets are provisional until Method is fully frozen.
%
% =====================================================================

\section{Introduction}
\label{sec:intro}

Graph-based retrieval-augmented generation, or GraphRAG, has become a
natural approach for multi-hop question answering because the evidence
needed to answer a question is often distributed across several
passages. Instead of ranking each passage independently, GraphRAG methods
organize a corpus into a structure of entities, relations, propositions,
passages, or summaries, and then use this structure to retrieve evidence
that is connected across documents. This direction is especially relevant
for multi-hop QA, where a useful retrieval result should not only contain
passages that resemble the question, but also expose bridge facts and
tail evidence that allow a reader to connect the question entities to the
answer.

Existing GraphRAG systems have developed several routes for introducing
structure into retrieval. One route builds an OpenIE-derived graph over
the corpus and propagates relevance from question-side seeds. HippoRAG
and HippoRAG 2 \citep{jimenez2024hipporag,gutierrez2025hipporag}, for
example, extract OpenIE facts, construct passage--phrase graphs, and use
personalized PageRank to spread relevance over a corpus-wide graph. This
route gives dense or lexical retrieval an explicit connectivity layer,
but its graph operator is still primarily a relevance propagator. In
practice, propagated mass can concentrate around the closest seed
neighborhood, while bridge passages several hops away receive weaker
scores. Moreover, when relation-bearing facts are converted into graph
connectivity, the final passage ranking still needs a separate mechanism
to decide whether the selected top-$K$ passages cover the different
entities and relations required by the question.

A second route makes multi-hop structure more explicit by searching or
constructing paths. PropRAG \citep{wang2025proprag} uses propositions as
context-rich graph units and applies beam search to discover reasoning
paths, while KG-oriented systems such as ToG \citep{sun2024tog} and
Relink \citep{huang2026relink} traverse or repair relation graphs at
query time. A third route changes the graph representation itself:
GraphRAG \citep{edge2024graphrag} retrieves through community summaries,
and hypergraph-based methods such as HGRAG \citep{hgrag2024} connect
entities and passages through higher-order structures. These approaches
address important weaknesses of flat passage ranking, but they also
leave a tension for retrieval-first multi-hop QA. Path-oriented methods
can make chains explicit, yet often rely on query-time search or
language-model-guided construction. Summary and hypergraph methods
improve corpus-level organization, yet their coarse or mixed-granularity
representations do not directly specify how propagated relevance should
be converted into a short evidence list for the reader.

This paper is motivated by the following gap. Graph structure helps
identify connected evidence, but GraphRAG systems ultimately have to
return a ranked passage list under a small reader budget. The problem is
therefore not only how to propagate relevance over a graph, but how to
turn propagated relevance into evidence that is both connected and
covering. It should be connected enough to preserve bridge relations
across documents, and covering enough to expose the distinct
question-side entities and relations that the reader must resolve. We
focus on this conversion problem while keeping the retrieval interface
simple. The system outputs a ranked evidence list $R_q$ of natural
language passages, rather than a symbolic path or a generated evidence
chain, and the downstream reader consumes a prefix of this list.

We propose \methodname{}, a graph-based retrieval framework built around
index-grounded evidence propagation. The key idea is to keep relation
structure on the indexing side and use query-time computation to
propagate, expand, and select evidence over a document-grounded OpenIE
graph. Offline indexing constructs a multi-layer graph that couples
passages, entity phrases, OpenIE facts, and provenance
(\S\ref{sec:method:substrate}). At query time, \methodname{} induces a
query-local view of this graph and runs fact-level personalized PageRank
to obtain a query-conditioned evidence flow $\pi_q$
(\S\ref{sec:method:retrieval}). Passage scores are read out from the
propagated masses of provenance-linked fact units. A retrieval
enhancement stage then follows relation-grounded edges to admit bridge
and tail-entity passages, and reranks candidates with a noisy-OR coverage
model over question-side requirements
(\S\ref{sec:method:enhancement}). In this way, \methodname{} treats graph
propagation as the start of multi-hop evidence selection rather than as
the final ranking by itself.

We evaluate \methodname{} on multi-hop QA benchmarks under a common
reader protocol, comparing against dense, graph-based, and
chain-oriented retrieval baselines. The evaluation measures retrieval
quality and downstream QA performance, and includes ablations for local
evidence propagation, relation-grounded expansion, and coverage-aware
reranking. These experiments are designed to test whether the proposed
conversion from graph relevance to ranked evidence improves both support
coverage and answer quality.
% TODO: Add final headline numbers once all main-table results are frozen.

% TODO: Finalize the three contribution bullets after the Method section is
% frozen. Current provisional wording follows.
The main contributions are as follows.
\begin{itemize}
    \item We introduce \methodname{}, a GraphRAG retrieval framework that
    models multi-hop evidence selection as query-conditioned propagation
    over document-grounded OpenIE facts and provenance-linked passages.
    \item We develop a retrieval mechanism that combines query-local
    fact-level personalized PageRank with relation-grounded expansion and
    noisy-OR coverage reranking over question-side requirements.
    \item We provide an empirical study that evaluates retrieval recall,
    downstream QA performance, and component ablations under a unified
    reader setting.
\end{itemize}
